\begin{frame}{Ví dụ minh họa: Phân rã thành các Luật Cục bộ (Local Rules)}
  Thay vì học bảng cộng cho toàn bộ chuỗi bit (ví dụ $011 + 101$), thuật toán được phân rã thành các \textbf{thao tác cục bộ} tại mỗi vị trí $i$:

  \begin{itemize}
    \item \textbf{Nhóm 1: Bitwise Addition (Tại bit $i$)}
      \begin{itemize}
        \item Quy tắc $A_i$: $(p_i=1, q_i=0) \to p_i=1$
        \item Quy tắc $B_i$: $(p_i=1, q_i=1) \to p_i=0, \text{sinh bit nhớ } c_i=1$
      \end{itemize}
    
    \item \textbf{Nhóm 2: Carry Propagation (Tại bit $i$)}
      \begin{itemize}
        \item Quy tắc $C_i$: Nếu có bit nhớ $c_i=1 \to$ cộng dồn vào $q_{i+1}$ (hoặc lan truyền tiếp).
      \end{itemize}
  \end{itemize}

  \vspace{0.2cm}
  $\Rightarrow$ Bài toán quy về việc học một tập hợp \textbf{"bảng chữ cái"} các bước chuyển đổi nhỏ.
\end{frame}

\begin{frame}{Chiến lược Encode 1: "Bag of Rules" & Training Data}
  \textbf{Tổng hợp Templates:}
  Với chuỗi $L$ bit, ta có tập hợp tất cả các quy tắc cục bộ cho mọi $i$:
  \[ T_{total} = \bigcup_{i=0}^{L-1} \{ \text{Rules tại } i \} \]

  \textbf{Mã hóa (Encoding):}
  Ta gộp tất cả $N$ rules cục bộ vào một không gian vector chung $\mathbb{R}^N$. Mỗi rule $t_k \in T_{total}$ được gán một chiều riêng biệt (one-hot).

  \textbf{Training Data (Identity Matrix):}
  \[
    X_{train} = I_N = 
    \begin{bmatrix}
    e_1 & e_2 & \dots & e_N
    \end{bmatrix}
  \]
  $\Rightarrow$ Mạng học song song tất cả các quy tắc cục bộ dưới dạng các vector \textbf{trực giao}.
\end{frame}

\begin{frame}{Chiến lược Encode 2: Inference & Lookup}
  \textbf{Câu hỏi:} Khi tính toán thực tế, mạng làm gì?

  \textbf{Inference (Superposition):}
  Tại một thời điểm, trạng thái hệ thống khớp với rule $t_k$. 
  \begin{itemize}
    \item Input vector $x$ chính là $e_k$.
    \item Nhờ tính chất trực giao của NTK trên $X_{train}$, output chỉ phụ thuộc vào label của $t_k$:
  \end{itemize}
  \[
    y(x) \approx \sum_{j} \underbrace{\langle e_k, e_j \rangle}_{\delta_{kj}} y^{(j)}_{train} = y^{(k)}_{train}
  \]

  \vspace{0.2cm}
  \textbf{Ý nghĩa:} Mạng không cần nhìn thấy "bức tranh toàn cảnh" của phép cộng, mà chỉ cần nhận diện đúng "mảnh ghép" (local rule) đang xuất hiện và thực thi nó.
\end{frame}
